两条平行水平理想导轨间距 \(\ell \)，右端用电阻 R>0 连接，一根质量 m、长 \(\ell \) 的导电棒在导轨上无摩擦滑动。装置位于垂直导轨平面向上的均匀磁场 B>0。电路电感、棒和导轨自身电阻均忽略，棒初始沿 +x 以速度 \(v_{0}>0\) 运动，之后没有外界机械推力。
(1) 由动生电动势和回路电流方向求棒受磁力；(2) 解出棒的速度、位移、电流随时间的完整表达式；(3) 求总电阻热并与初始动能比较；(4) 研究最终位移以及电阻趋于无穷和磁场消失的奇异极限，解释为何电磁能流与机械耗散相容。
